\section{The Redshift of Distant Galaxies}
\label{sec:redshift}

This section examines the cosmological redshift phenomenon and its implications for understanding the expanding universe within the I-theory framework.

\subsection{Cosmological Redshift}
\label{subsec:cosmological_redshift}

The redshift of light from distant galaxies is quantified by:

\begin{equation}
\label{eq:redshift_definition}
z = \frac{\lambda_{\text{observed}} - \lambda_{\text{emitted}}}{\lambda_{\text{emitted}}} = \frac{\Delta\lambda}{\lambda_{\text{emitted}}}
\end{equation}

where:
\begin{itemize}
    \item $\lambda_{\text{observed}}$ is the wavelength of light received on Earth
    \item $\lambda_{\text{emitted}}$ is the original wavelength emitted by the galaxy
    \item $z$ is the dimensionless redshift parameter
\end{itemize}

\subsection{Physical Origin of Redshift}

The redshift arises due to the expansion of the universe, stretching the wavelengths of light as they travel across cosmic distances.

\subsubsection{Scale Factor Relation}

The redshift is directly related to the cosmic scale factor:

\begin{equation}
\label{eq:scale_factor_redshift}
1 + z = \frac{a(t_{\text{obs}})}{a(t_{\text{emit}})} = \frac{a_0}{a(t_e)}
\end{equation}

where:
\begin{itemize}
    \item $a(t)$ is the scale factor of the universe at time $t$
    \item $a_0 = a(t_{\text{obs}})$ is the present-day scale factor (conventionally set to 1)
    \item $t_e$ is the emission time
\end{itemize}

\subsubsection{Hubble's Law}

For nearby galaxies ($z \ll 1$), redshift is approximately proportional to distance:

\begin{equation}
\label{eq:hubble_law}
cz \approx H_0 d
\end{equation}

where:
\begin{itemize}
    \item $c$ is the speed of light
    \item $H_0$ is the Hubble constant ($H_0 \approx 67.4$ km s$^{-1}$ Mpc$^{-1}$)
    \item $d$ is the proper distance to the galaxy
\end{itemize}

\subsection{High Redshift Regime}

At high redshifts ($z > 6$), we observe light from the era when the universe was only a few hundred million years old.

\subsubsection{Lookback Time}

The lookback time as a function of redshift:

\begin{equation}
\label{eq:lookback_time}
t_L(z) = t_0 - t(z) = \int_0^z \frac{dz'}{(1+z')H(z')}
\end{equation}

where $t_0$ is the age of the universe and:

\begin{equation}
\label{eq:hubble_parameter}
H(z) = H_0 \sqrt{\Omega_m(1+z)^3 + \Omega_r(1+z)^4 + \Omega_\Lambda}
\end{equation}

with:
\begin{itemize}
    \item $\Omega_m$ = matter density parameter $\approx 0.315$
    \item $\Omega_r$ = radiation density parameter $\approx 9.2 \times 10^{-5}$
    \item $\Omega_\Lambda$ = dark energy density parameter $\approx 0.685$
\end{itemize}

\subsubsection{Comoving Distance}

The comoving distance to an object at redshift $z$:

\begin{equation}
\label{eq:comoving_distance}
d_c(z) = \frac{c}{H_0} \int_0^z \frac{dz'}{\sqrt{\Omega_m(1+z')^3 + \Omega_r(1+z')^4 + \Omega_\Lambda}}
\end{equation}

\subsection{I-Theory Modifications}

In I-theory, the expansion rate is modified by I-field contributions:

\begin{equation}
\label{eq:i_modified_hubble}
H_I(z) = H(z) \sqrt{1 + \frac{\Omega_I}{I}(1+z)^{3(1+w_I)}}
\end{equation}

where:
\begin{itemize}
    \item $\Omega_I$ is the I-field density parameter
    \item $w_I$ is the I-field equation of state parameter
\end{itemize}

\subsubsection{Modified Redshift Relation}

The I-modified redshift-distance relation:

\begin{equation}
\label{eq:i_redshift_distance}
d_L(z) = (1+z) d_c(z) \left(1 + \alpha_I \frac{z^2}{1+z}\right)
\end{equation}

where $\alpha_I$ is the I-correction factor and $d_L$ is the luminosity distance.

\subsection{Observational Signatures}

\subsubsection{Type Ia Supernovae}

The distance modulus for Type Ia supernovae:

\begin{equation}
\label{eq:distance_modulus}
\mu(z) = 5\log_{10}\left[\frac{d_L(z)}{\text{Mpc}}\right] + 25
\end{equation}

I-theory predicts deviations from $\Lambda$CDM:

\begin{equation}
\label{eq:i_deviation}
\Delta\mu_I(z) = \mu_I(z) - \mu_{\Lambda\text{CDM}}(z) \approx 5\alpha_I \frac{z^2}{(1+z)\ln 10}
\end{equation}

\subsubsection{Baryon Acoustic Oscillations}

The BAO scale at redshift $z$:

\begin{equation}
\label{eq:bao_scale}
r_s(z_d) = \int_{z_d}^\infty \frac{c_s(z)}{H(z)} dz
\end{equation}

where $c_s$ is the sound speed and $z_d$ is the drag epoch redshift.

I-modifications alter the angular diameter distance:

\begin{equation}
\label{eq:angular_diameter}
d_A^I(z) = \frac{d_c^I(z)}{1+z} = d_A^{\Lambda\text{CDM}}(z) \left(1 - \beta_I \frac{z}{1+z}\right)
\end{equation}

\subsection{Extreme Redshift Observations}

\subsubsection{High-z Galaxies}

For $z > 10$ galaxies:

\begin{equation}
\label{eq:high_z_age}
t_{\text{universe}}(z) = \frac{1}{H_0} \int_z^\infty \frac{dz'}{(1+z')\sqrt{\Omega_m(1+z')^3 + \Omega_\Lambda}}
\end{equation}

At $z = 10$: $t_{\text{universe}} \approx 480$ Myr.

\subsubsection{Lyman-Alpha Forest}

The Ly$\alpha$ forest probes the intergalactic medium:

\begin{equation}
\label{eq:lya_optical_depth}
\tau_{\text{eff}}(z) = \int_0^{z_{\text{QSO}}} \sigma_{\alpha} n_{HI}(z') \frac{c\,dz'}{H(z')}
\end{equation}

I-theory predicts modifications:

\begin{equation}
\label{eq:i_optical_depth}
\tau_I(z) = \tau_{\text{eff}}(z) \exp\left(\gamma_I \frac{z}{1+z}\right)
\end{equation}

\subsection{Cosmic Microwave Background}

The last scattering surface occurs at $z \approx 1090$:

\begin{equation}
\label{eq:cmb_redshift}
T_{\text{CMB}}(z) = T_0 (1+z) = 2.725\, \text{K} \times (1+z)
\end{equation}

At last scattering: $T_{\text{CMB}}(z_{ls}) \approx 2970$ K.

\subsubsection{Temperature Anisotropies}

The CMB angular power spectrum:

\begin{equation}
\label{eq:cmb_power_spectrum}
C_\ell = \frac{2}{\pi} \int dk\, k^2 \mathcal{P}(k) |\Theta_\ell(k)|^2
\end{equation}

where $\Theta_\ell(k)$ is the radiation transfer function.

I-corrections modify the acoustic peaks:

\begin{equation}
\label{eq:i_acoustic_peaks}
\ell_n^I = \ell_n^{\text{std}} \left(1 + \epsilon_I \frac{n^2}{\ell_n}\right)
\end{equation}

\subsection{Redshift Evolution of Physical Quantities}

\subsubsection{Temperature Evolution}

For photons:
\begin{equation}
\label{eq:photon_temperature}
T_\gamma(z) = T_0(1+z)
\end{equation}

For matter (adiabatic expansion):
\begin{equation}
\label{eq:matter_temperature}
T_m(z) \propto (1+z)^2
\end{equation}

\subsubsection{Density Evolution}

Matter density:
\begin{equation}
\label{eq:matter_density}
\rho_m(z) = \rho_{m,0}(1+z)^3
\end{equation}

Radiation density:
\begin{equation}
\label{eq:radiation_density}
\rho_r(z) = \rho_{r,0}(1+z)^4
\end{equation}

I-field density:
\begin{equation}
\label{eq:i_field_density}
\rho_I(z) = \rho_{I,0}(1+z)^{3(1+w_I)} \otimes I
\end{equation}
